\documentclass[10pt,a4paper]{article}
\usepackage{fontspec}
\setmainfont{Liberation Sans}
\usepackage{polyglossia}
\setdefaultlanguage{polish}
\usepackage[a4paper,margin=19mm,headheight=15pt]{geometry}
\usepackage{xcolor,tabularx,booktabs,array,enumitem,fancyhdr,titlesec,hyperref}
\definecolor{navy}{HTML}{263650}
\definecolor{blue}{HTML}{5475A5}
\definecolor{pale}{HTML}{EFF3F8}
\definecolor{muted}{HTML}{5D6878}
\hypersetup{colorlinks=true,urlcolor=blue}
\pagestyle{fancy}
\fancyhf{}
\fancyhead[L]{\small\color{muted}Tutorland}
\fancyhead[R]{\small\color{muted}Mapa stron}
\fancyfoot[C]{\small\color{muted}\thepage}
\renewcommand{\headrulewidth}{0.3pt}
\setlength{\parindent}{0pt}
\setlength{\parskip}{5pt}
\setlength{\emergencystretch}{3em}
\sloppy
\titleformat{\section}{\Large\bfseries\color{navy}}{\thesection.}{0.55em}{}
\newcolumntype{Y}{>{\raggedright\arraybackslash}X}
\begin{document}
\begin{titlepage}
\thispagestyle{empty}
\vspace*{32mm}
{\large\color{blue}\textbf{DOCELOWA STRUKTURA PRODUKTU}}
\vspace{12mm}
{\Huge\bfseries\color{navy}Tutorland\par}
\vspace{4mm}
{\LARGE Wszystkie potrzebne strony i widoki\par}
\vspace{13mm}
\color{blue}\rule{\textwidth}{2pt}
\vspace{9mm}
{\large Mapa ekranów publicznych i widoków kont ucznia oraz korepetytora.}
\vfill
\colorbox{pale}{\parbox{\dimexpr\textwidth-2\fboxsep\relax}{\textbf{Jak czytać dokument}\\
Mapa opisuje aktualne widoki aplikacji. Część formularzy i podsekcji jest stanem jednej strony; trasy są obsługiwane przez frontend React.}}\par
\vspace{12mm}
{\small\color{muted}Struktura oparta na zakresie produktu Tutorland: uczniowie i korepetytorzy, rezerwacje 40-minutowe, Zoom, spotkania i opinie.}
\end{titlepage}

\section{Strony publiczne}
\small
\begin{tabularx}{\textwidth}{>{\bfseries}p{8mm}p{41mm}p{30mm}Y}
\toprule
Nr & Strona / przykładowa trasa & Dostęp & Zawartość i działanie\\
\midrule
1 & Strona główna i katalog\newline \url{/} & Każdy & Karty profili, wyszukiwanie po imieniu i nazwisku, filtry przedmiotów i poziomów oraz stany ładowania i braku wyników. Uczeń pozostaje w katalogu po zalogowaniu.\\
2 & Profil korepetytora\newline \url{/?tutor=:tutorId} & Każdy & Opis, przedmioty, poziomy, cena, oceny, opinie i dostępne terminy. Gość może przeglądać profil i terminy, ale rezerwacja wymaga konta ucznia.\\
\bottomrule
\end{tabularx}
\normalsize

\section{Konto i wspólne widoki prywatne}
\small
\begin{tabularx}{\textwidth}{>{\bfseries}p{8mm}p{41mm}p{30mm}Y}
\toprule
Nr & Strona / przykładowa trasa & Dostęp & Zawartość i działanie\\
\midrule
3 & Rejestracja\newline \url{/register} & Gość & Utworzenie konta z wyborem roli ucznia lub korepetytora oraz walidacją danych.\\
4 & Logowanie\newline \url{/login} & Gość & Logowanie z wyborem roli konta. Po logowaniu użytkownik trafia do odpowiedniego panelu lub wraca do strony profilu.\\
5 & Spotkania i powiadomienia\newline \url{/notifications} & Uczeń / korepetytor & Własne nadchodzące i zakończone spotkania, status oraz link Zoom. Uczeń zarządza opiniami, korepetytor widzi rezerwacje i nowe opinie.\\
6 & Panel profilu i dostępności\newline \url{/profile} & Korepetytor & Edycja własnego profilu, ceny, przedmiotów i poziomów; tygodniowe godziny przez suwaki, kalendarz rezerwacji i dni wolne na kolejny miesiąc.\\
7 & Opinia\newline Formularz w \url{/notifications} & Uczeń po lekcji & Ocena 1–5 oraz komentarz. Jedna aktywna opinia na korepetytora; można ją edytować, usunąć i po usunięciu dodać ponownie.\\
8 & Połączenie z Zoom\newline \url{/api/zoom/connect} i \url{/api/zoom/callback} & Korepetytor & OAuth konta Zoom, pobranie dostępnego zdjęcia profilowego i utworzenie 40-minutowych spotkań. Wymaga konfiguracji aplikacji Zoom na serwerze.\\
\bottomrule
\end{tabularx}
\normalsize
\end{document}
